\maketitle

\begin{abstract}
We study oracle classification error along a fixed prefix filtration and keep
the scope deliberately conditional: once the filtration and reference measure
are fixed, the Bayes risk $\varepsilon_m(\mathsf{P};\mu)$ decays exactly as
the weighted boundary mass.  The structural input is the discrete Tanaka
identity $\varepsilon_m=\varepsilon_0-\EE[\mathcal{T}_m]$ applied to the posterior
martingale, together with a weighted boundary law expressing
$\varepsilon_m$ as the boundary-mass sum of local ambiguities.  Standard
count-based upper and lower bounds follow whenever cylinders of the same
length are uniformly comparable, and the same decomposition enforces lower
growth of mixed cylinders from any prescribed lower bound on $\varepsilon_m$.
The genuinely new content is an intrinsic hidden-Markov/sofic theorem and an
exact finite-state scan-error refinement:
for primitive absorbing prefix filters on mixing shifts of finite type with the
Parry measure, the mixed posterior states form a finite strongly connected
graph whose spectral radius simultaneously yields automatic thickness and
determines the exponential growth of boundary cylinders.  The finite posterior
state space also gives an exact transfer-matrix formula for the scan error,
including a rational generating function and an explicit Perron asymptotic
constant.  This places the
conditional rate laws inside the classical hidden-Markov literature while
remaining honest about scope.  The worked examples include a
positive-measure run-length hidden-Markov example, the parity and self-similar
cases, and a boundary-growth failure illustrating the necessity of
thickness.
A polynomial-distortion variant is recorded separately as an auxiliary
remark rather than as a Gibbs theorem.
\end{abstract}

\tableofcontents

\section{Introduction}
\label{sec:intro}

This note studies a fixed-filtration oracle question at the level appropriate
for a short communication: how much does refining prefix information improve
the best possible decision error for a target event?  Let
$(X,\mathcal{B},\mu)$ be a probability space with an increasing filtration
$(\mathcal{F}_m)_{m\ge 0}$ generated by $\{0,1\}$-valued functions
$s_0,s_1,\ldots$  The scan error profile
\[
m\longmapsto\varepsilon_m(\mathsf{P};\mu)
:=\inf\bigl\{\mu(\mathsf{P}\triangle\mathsf{Q}):\mathsf{Q}\in\mathcal{F}_m\bigr\}
\]
records the smallest classification error at each prefix depth.

The Bayes cylinder decomposition and the optimality of the threshold classifier
$p_m\ge 1/2$ are classical
\cite{Tsybakov2009Nonparametric,DevroyeGyorfiLugosi1996}.  We recall them in
Section~\ref{sec:scan} only in the form needed later.  The only sources of
genuinely new mathematics are the hidden-Markov/sofic theorem and the exact
finite-state scan-error transfer theorem stated below;
all remaining statements are reorganizations of standard identities inside a
fixed filtration.  The main points are:
\begin{enumerate}[label=(\roman*)]
\item A \emph{discrete Tanaka identity}
  (Proposition~\ref{prop:scan-tanaka}) connecting error decay to the accumulation
  of threshold crossings of the martingale $(p_m)$.  This is a structural
  observation that follows from the standard discrete-time Tanaka--Meyer
  formula \cite{LochowskiOblojPromelSiopaes2021}.
\item A \emph{weighted boundary law}
  (Proposition~\ref{prop:weighted-boundary-law}) expressing $\varepsilon_m$ as
  boundary mass times averaged local ambiguity, together with count-based
  consequences and a converse-style lower-growth statement
  (Propositions~\ref{prop:boundary-dimension}--\ref{prop:volume-boundary-scaling}
  and Proposition~\ref{prop:boundary-growth-necessary}).
\item An \emph{intrinsic hidden-Markov/sofic class}
  (Definition~\ref{def:hidden-markov-filter} and
  Theorem~\ref{thm:hidden-markov-thickness}) whose mixed posterior states form
  a finite strongly connected graph.  Primitive absorbing prefix filters thereby
  deliver automatic thickness and a spectral formula for the boundary growth,
  while Theorem~\ref{thm:exact-scan-error-transfer} gives the exact
  transfer-matrix formula, rational generating function, and leading Perron
  constant for the scan error itself,
  placing the conditional rate laws squarely inside the standard literature on
  sofic factors and probabilistic automata
  \cite{Blackwell1957,Heller1965,BaumPetrie1966,BoylePetersenHMP,PollicottKempton2016,Verbitskiy2011,Rabin1963,Paz1971,PerrinPin2004}.
\item \emph{Worked examples and failure modes}: a benchmark hidden-Markov
  factor example, the first-hit parity event, the self-similar boundary exponent
  $\frac12$, and failure examples showing the necessity of thickness, the effect
  of filtration choice, and the detailed computation for
  Example~\ref{ex:thickness-failure}.
\end{enumerate}
The novelty claim is therefore precise: the hidden-Markov/sofic theorem and
the exact finite-state scan-error transfer theorem are the new statements,
whereas the remaining material is recorded so that the conditional rate laws
can be cited cleanly.  The structural observations are
nonetheless useful because they isolate what is genuinely filtration-driven
and what must be assumed from outside (thickness, cylinder comparability).
Section~\ref{sub:related} provides pointers to the closest partition, hidden
Markov, and Gibbs literatures.
To keep the focus on Theorems~\ref{thm:hidden-markov-thickness}
and~\ref{thm:exact-scan-error-transfer}, the
preliminary Sections~\ref{sec:prefix}--\ref{sec:tanaka} and the auxiliary
log-Gibbs discussion in Section~\ref{sub:nonuniform} have been shortened to
bare statements with explicit citations; routine proofs are omitted.

\subsection{Heuristic comparison with margin assumptions in classification}
\label{sub:margin}

The thickness hypothesis in Proposition~\ref{prop:volume-boundary-scaling}
should be read only as a heuristic comparison with margin and noise conditions
in nonparametric classification.
Mammen and Tsybakov \cite{MammenTsybakov1999} introduce a margin condition
with exponent $\alpha\ge 0$: for the regression function
$p_\infty(x)=\mu(\mathsf{P}\mid\mathcal{F}_\infty)(x)$,
\[
\mu\!\left(\left|p_\infty(x)-\tfrac12\right|\le t\right)\le C_0\,t^\alpha
\qquad\text{for all }t>0.
\]
Tsybakov \cite{Tsybakov2004} extends this to a general noise-exponent
framework, showing that margin conditions control the rate of optimal
plug-in classifiers.

Our uniform thickness condition
$\min\{\mu(\mathsf{P}\mid C),1-\mu(\mathsf{P}\mid C)\}\ge\theta>0$ for all
boundary cylinders $C\in\partial_m\mathsf{P}$ is a \emph{uniform two-sided
nondegeneracy} assumption on individual mixed cylinders: it demands that
every boundary cylinder carries a bounded-away-from-zero fraction of
both $\mathsf{P}$ and $\mathsf{P}^c$.  The Mammen--Tsybakov condition is,
by contrast, a \emph{global distributional} statement about the regression
function $p_\infty$ near $\frac12$; the exponent $\alpha$ is not a fixed
threshold but a rate parameter.

The two conditions are related but structurally different:
\begin{enumerate}[label=(\arabic*)]
\item Uniform thickness prevents one-sided concentration on individual
  cylinders, while the margin condition integrates over the whole space.
  A measure satisfying the margin condition with large $\alpha$
  can still have individual boundary cylinders with $\mu(\mathsf{P}\mid C)\to 0$,
  violating uniform thickness.
\item Thickness controls \emph{oracle} error (best decision with known measure),
  whereas the margin exponent controls \emph{estimation} error (rate of
  learning from data).  These are different quantities, and the fixed
  constant $\theta$ does not play the same role as $\alpha$.
\end{enumerate}
For a more direct point of comparison in the dyadic-partition context, Scott
and Nowak \cite{ScottNowak2006} combine a dyadic box-counting boundary class
with noise assumptions that are closer in spirit to our setup; see
Section~\ref{sub:related} for a detailed comparison.

\subsection{Related work}
\label{sub:related}

Bayes risk on partitions and optimality of threshold classifiers are classical;
see \cite{DevroyeGyorfiLugosi1996} and \cite{AudibertTsybakov2007} for refined
plug-in rates.  The underlying martingale convergence is Doob's upward theorem
\cite{Doob1953,Williams1991,Durrett2019}.  The discrete Tanaka formula is a
lattice analogue of Tanaka's formula for Brownian local time
\cite{Tanaka1963,RevuzYor1999}; Proposition~\ref{prop:scan-tanaka} is used here
only as a structural identity, not as a new Tanaka theorem.

\paragraph{Classical hidden-Markov references.}
The probabilistic-function program initiated by
\cite{Blackwell1957,Heller1965,BaumPetrie1966} already describes finite-state
filters of Markov chains, while \cite{Thomas1990,PerrinPin2004,Kitchens1998}
develop the automata-on-infinite-words and symbolic-dynamics background that we
invoke.  Section~\ref{sec:rates} uses this machinery only to the extent needed
to identify when the fixed-filtration oracle quantities inherit automatic
thickness and Perron--Frobenius growth.

Within that classical landscape, Theorem~\ref{thm:hidden-markov-thickness}
chooses the absorbing prefix-filter viewpoint: we keep the reference filtration
and measure fixed, treat the finite-state filter as a transducer producing the
oracle posterior, and ask only for quantitative consequences (automatic
thickness and explicit boundary growth) for that filtration.  Thus the result
does not reprove the existence of hidden-Markov factors in the sense of
\cite{Blackwell1957,BaumPetrie1966} or of omega-regular acceptance conditions
\`a la \cite{Thomas1990,PerrinPin2004}; instead it pinpoints how those classical
structures feed into the scan-error decay studied here.

Cylinder partition structure in symbolic dynamics is treated in
\cite{LindMarcus1995}; boundary combinatorics connect to
\cite{Fischer1975,Krieger1984,Boyle1983}.  Exponential cylinder estimates for
the Parry measure (measure of maximal entropy) are in \cite{Parry1964}; for
general Gibbs measures the corresponding cylinder bounds are
potential-dependent rather than uniform over all cylinders of a fixed length
\cite{Bowen1975,WaltersErgodicTheory}.
The Shannon--McMillan--Breiman theorem \cite{Shannon1948,McMillan1953,Breiman1957}
gives almost-sure asymptotics for the cylinder \emph{containing} a given
sample point; this is an almost-sure pointwise statement, weaker than the
uniform all-cylinder bounds assumed in our rate theorems (see
Section~\ref{sec:discussion}).  Background on dynamics and information theory
is in \cite{WaltersErgodicTheory,LindMarcus1995,CoverThomas2006,MorseHedlund1940}.

\paragraph{Closer neighboring statistical literature.}
A closer comparison point than nonlinear approximation alone is the literature
on adaptive partition classification and level-set estimation.  Binev et al.\
\cite{BinevCohenDahmenDeVore2014} study adaptive partition classifiers as set
approximators to the Bayes set.  Willett and Nowak
\cite{WillettNowak2007}, Singh et al.\ \cite{SinghScottNowak2009}, and Scott
and Davenport \cite{ScottDavenport2007} analyze level-set estimation with
hierarchical partitions, Hausdorff loss, or cost-sensitive reductions.
Historically, CART \cite{BreimanFriedmanOlshenStone1984} is the baseline tree
partition framework behind much of this literature.  The distinction from the
present paper is threefold: the measure is known rather than estimated, the
filtration is fixed rather than adaptively selected, and the quantity of
interest is the exact oracle Bayes risk along that single filtration rather
than minimax estimation risk.  In that sense the present results lie on the
decision-theoretic side of partition refinement; compare Blackwell's
comparison-of-experiments viewpoint \cite{Blackwell1951}.

\paragraph{Comparison with dyadic decision-tree and approximation literature.}
The boundary-count viewpoint is closest to dyadic decision-tree work in
statistical learning, especially Scott and Nowak \cite{ScottNowak2006} and
Blanchard et al.~\cite{BlanchardSchaferRozenholcMuller2007}.  Scott and
Nowak study \emph{sample-based} classification with dyadic decision trees
under a dyadic box-counting boundary class together with noise assumptions.
Blanchard et al.\ analyze optimal dyadic decision trees and oracle
inequalities for data-driven tree selection.  A second comparison point is
the adaptive-partition approximation literature of Donoho
\cite{Donoho1997}, DeVore \cite{DeVore1998}, and Cohen et al.\
\cite{CohenDahmenDaubechiesDeVore2001}, where dyadic trees encode functions
through best-basis or best-tree approximants.  Our setting is narrower: the
measure is known, the filtration is fixed in advance, and the object of study
is the \emph{oracle prefix-Bayes error} $\varepsilon_m(\mathsf{P};\mu)$ for a
single prescribed chain of prefix partitions.  The shared ingredient is
therefore the count of mixed dyadic cells, but the conclusions differ:
here we obtain direct oracle identities and rate bounds along that fixed
filtration, rather than minimax learning rates, empirical oracle
inequalities, or best-tree approximation theorems.

\begin{table}[ht]
\centering
{\scriptsize
\setlength{\tabcolsep}{3pt}
\begin{tabular}{@{}>{\raggedright\arraybackslash}p{1.00in}>{\raggedright\arraybackslash}p{0.70in}>{\raggedright\arraybackslash}p{0.76in}>{\raggedright\arraybackslash}p{0.82in}>{\raggedright\arraybackslash}p{1.18in}>{\raggedright\arraybackslash}p{1.00in}@{}}
\toprule
reference & setting & object & information & complexity / noise & conclusion \\
\midrule
Scott--Nowak \cite{ScottNowak2006} &
dyadic classification &
classifier / boundary &
sampled data &
box-counting boundary class plus noise assumptions &
minimax classification rates \\
Blanchard et al.\ \cite{BlanchardSchaferRozenholcMuller2007} &
empirical dyadic trees &
tree classifier &
sampled data &
tree penalty and oracle model selection &
oracle inequalities; optimal trees \\
Donoho \cite{Donoho1997}; DeVore \cite{DeVore1998}; Cohen et al.\ \cite{CohenDahmenDaubechiesDeVore2001} &
adaptive approximation &
function / encoding &
target known exactly &
best-basis or tree budget; no noise layer &
best-tree approximation and optimal encoding \\
this paper &
fixed prefix filtration &
event indicator / Bayes decision &
known measure &
mixed-cylinder count plus thickness on mixed cells &
oracle decay of $\varepsilon_m$ \\
\bottomrule
\end{tabular}
\par}
\caption{Side-by-side positioning against the closest dyadic-tree and adaptive-approximation references.}
\label{tab:dyadic-comparison}
\end{table}

\paragraph{Hidden-Markov and Gibbs viewpoints.}
Hidden-Markov and sofic measures are treated extensively in the symbolic
dynamics literature \cite{BoylePetersenHMP,Verbitskiy2011}, building on the
classical probabilistic-function program of
\cite{Blackwell1957,Heller1965,BaumPetrie1966}.  Automata on infinite objects
and symbolic-coding references such as
\cite{Thomas1990,PerrinPin2004,Kitchens1998} provide the complementary
background for the automaton formulation used here.  Factors of Gibbs measures
give the dynamical context for potential-weighted cylinder estimates
\cite{PollicottKempton2016}, and the probabilistic-automata models of Rabin and
Paz \cite{Rabin1963,Paz1971} already encode the finite-state filters that
appear in Section~\ref{sec:rates}.  The present paper does not attempt a full
thermodynamic treatment of hidden-Markov factors; our contribution is instead to
identify exactly where such classical structures provide automatic thickness and
spectral growth inside the fixed-filtration oracle picture.

